\documentclass[12pt]{article}

\usepackage[a4paper, total={6in, 8in}]{geometry}
\usepackage{textcomp}

\begin{document}
\setlength{\parindent}{0pt}

\def \t {\textrightarrow}
\def \v {\vspace{1em}}
\def \bi {\begin{itemize}}
\def \ei {\end{itemize}}
\def \s[#1] {\section*{#1}}
\def \ss[#1] {\subsection*{#1}}
\def \sss[#1] {\subsubsection*{#1}}

\s[Guerra fredda]
\ss[Afghanistan]
Rivoluzione in Iran, Komeini al potere \t guerra Iran e Iraq, ma c'è un altro fronte, quello in Afghanistan \\
La sconfitta in Vietnam aveva indebolito iniziativa USA \t quindi URSS aumenta influenza e armamenti \\
L URSS impone poi con una guerra rovinosa il suo controllo sull'Afghanistan \t nel 1978 era andato al potere un regime comunista filosovietico anche se grande opposizione di momenti religiosi e gruppi etnici \\
Macato appoggio popolare rende regime filosovietico instabile \t quindi nel 79 le truppe sovietiche entrano in Afghanistan con lo scopo di difendere uno stato alleato \\
Breznev è al potere ora, che era succeduto a Cruschev \\
I russi fanno grande fatica anche se hanno mezzi e appoggio dell'esercito regolare afagano \t ma il malcontento si è organizzato in guerriglia e viene supportato da USA e Pakistan \\
Guerrglia si nasconde sulle montagne e attacca le città \t ci saranno milioni di profughi che vanno in Pakistan \\
I guerriglieri, che erano gruppi religiosi e etnici, vengono sostenuti dagli USA \\
L'occupazione dell'Afghanistan sembra essere lo stesso evento della guerra del Vietnam ma per urss \\
Nel 87 Gorbachov annuncia ritiro delle truppe da Afghanistan \t l'impegno militare è stato di anni \\
Si aprogono negoziati a Ginvera tra USA, URSS, Af. e Pakistan \t e si arriva nell 88 a un accordo che prevede allontanamento delle truppe sovietiche da territorio \\
Il regime comunista in afg. pero senza urss cade per mano dei talebani, che salgono al potere del 92

\v

Negli anni 80 sale al potere Gorbachov \t mentre in USA c'è Ronald Regan, che è un attore di Hollywood e repubblicano \\
Loro due saranno i fautori della fine della guerra fredda, ma la strada sara difficile \t Regan sara all inizio una specie di Truman e sara rigido \\
Dichiara volonta di rispondere a qualsiasi attacco sovietico in modo armato \\
L'URSS porta avanti rafforzamento dell'apparato militare \t aleggia "la paralisi nucleare" \t se tutti e due hanno l'arma atomica, non c'è il + forte e il rischio e distruzione totale \\
Formulata la strategia MUD \t entrmabe potenze dovevano mantenere energia nucleare, equilibrio nucelare stabile \t nel caso di un attaccato dall'altro, l'altro distrugge allo stesso tempo l'altro \t chiunque da via a guerra nucleare viene distrutto, per cui nessuno la inizierà mai \\
In realta c'erano stati incontri precedenti \t nel 69 ci sono stati Colloqui SALT tra Breznev e Nixon \t dovevano limitare reciprocamente le testate nucleari \\
Allo stesso tempo conferenza Helsinki \t usa, urss e alri 34 paesi \t tema è sicurezza e cooperazione in Europa \t tutti paesi partecipanti si impegnano a garantire il rispetto dei diritti fondamentali e libertà politiche fondamentali

\v

All'inizio degli anni 80 pero Regan chiude \t per esempio potenzia apparato difensivo usa, progettando lo scudo spaziale = sistema difensivo pensato in funzione antisovietica \t Regan chiama URSS l'impero del male \\
Arriva pero Gorbachov, che vuole cambiare \t lui non condivide nulla della politica sovietica fino a quel momento, parla quindi di Glasnov (trasparenza) e Perestroika \t per cui sconfessa la politica chiusa e poco trasparente dell urss, con critiche interne \\
Si dichiara aperto a dialogo con Regan \t e lui nonostante tutto si mostra disponibile ad ascoltare \\
Ci saranno colloqui internazionali, 85 a Ginevra e 86 a Rekavik \t che portera i due a firmare accordo a ritare tutte testate nucleari in europa, nell 86 \\
Nell 88 Gorbachov fa annuncio all'ONU che l URSS ridurra in modo unilaterale la presenza delle truppe nell est europeo \t mezzo milione di uomini tornano in russia dai paesi europei \\
Il passo successivo è il crollo del muro \t sta dicendo che l urss fa passo in dietro \\
Nell 89 il muro di berlino viene distrutto dai cittadini di berlino est \\
Nel 91 crolla l URSS \t e nasce la CSI \\
Nel 91 si conclude anche accordo START \t distruzione del 25 percento dlele testate nucleari possedute dalle due potenze \t si delinea nuova epoca, chiamata "del terzo dopoguerra"

\v

Con il crollo dell urss la guerra fredda finisce \t ma questo non vuol dire pace \t il mondo continua ad essere attraversato da conflitti, con pero modalita diversa \\
Cio che finisce con guerra fredda è il bipolarismo mondiale, che pero aveva garantito ordine \t inizio anni 90 si presenta un disordine internazionale \\
Ci sono state una serie di guerre che non hanno garantito piu un periodo di pace mondiale \t questi conflitti non sono riconducibile a un nessun ordine mondiale \\
Questo è l'"epoca del disordine mondiale" \t non ci sono guerre dettate da uno schieramento internazionale \t prima c'era grande tensione, pero all'interno del bipolarismo \\
Il 21 secolo si apre con una guerra con l'attacco dell torri gemelle, 3000 morte \t primo attacco diretto in USA \\
Le guerre dopo 9:11 sono sempre state definite "guerre assimetriche" (termine introdotto da Andre Mack nel 75, con cui intende che ci sono guerre combattute tra forze che non sono confrontabili militarmente, strategicamente, etc.) \t es. Vietnam che ha fatto arretrare USA \\
Anche Isreale e palestina \t esercito organizzato contro guerriglieri, che pero continua \t sono guerre squilibrate il cui esito non è mai pero scontato \\
Durante le torri gemelle era Bush che aveva preso il posto di Clinton nel 2000 \t e da qua inizia guerra in Iraq e Afghanistan

\s[In Italia]
12 dicembre del 69 scoppia bomba in piazza fontana \t inizia con questo attentato il periodo del terrorismo politico che attraversa l italia negli anni 70 \\
Si parla di terrorismo di sinistra e di destra \t i terroristi di destra si sentono eredi ella repubblica di Salo \t l'attacco allo stato è per fare cadere repubblica e far riportare governo fascista \\
Terr. di sinistra portano avanti la resistenza e sono eredi dei partigiani \t la resistenza era stata tradita dal partito comunista, che aveva rinunciato alla rivoluzione ed era sceso a patto con valori borghesi \\
Terrorismo di sinistra è continuamento del 68 e dell'autunno caldo o non c'è continuità? dibattuto \t alcuni dicono che:
\bi
  \item no perche questo terrorismo ha organizzazione militaristica e gerarchica, contraria al sentimento di liberazione e di lotta del 68
  \item si perche anche le manifestazioni studentsche e autunno caldo abbia mostrato episodi di violenza, in cui c'erano bombe moltov e si predicava lotta armata
\ei
Queste pero sono interpretazioni \\
I due terrorismi hanno modalita:
\bi
\item stragista quello di destra \t episodio eclatante con morti e feriti
\item violenza ad personam di sinistra \t colpito lo stato nei suoi rappresentanti \t colpisco la realta dello stato colpendo person precise
\ei
Il terrorismo nero è il protagonista dei primi anni 70 \t all'inizio ci sono stragi \t che hanno anche l'intento di seminare panico tra la gente, e rendere impossibile vita quotidiana \t rendo difficile vita sociale per fare entrare in crisi \\
A brescia nel 74 esplode bomba, treno Italicus nel 74 in provincia di Bologna \t esplode bomba sul treno, con 12 morti e 48 feriti \t 1980 stazione di bologna, 85 morti e 200 feriti \\
Questo considerato attentato + sanguinoso del dopoguerra \\
Questo in parallelo a conflitti sociali, soprattutto al sud \t es. reggio calabria??? \t grande instabilita, scoperte anche collusioni di servizi segreti con terrorismo nero \\
Si vuole quindi creare clima per colpo di stato 
\end{document}
